\documentclass{article}
\usepackage{amsmath}
\usepackage{amssymb}
\usepackage{geometry}
\geometry{margin=1in}

\title{Aether Sparse-Event Kernel Notes}
\author{Aether Lang}
\date{}

\begin{document}
\maketitle

\section*{Status}

This file is retained as an archived research-note entrypoint. It is not the
active source of public claims for Aether Lang.

The active documentation standard is maintained in:

\begin{itemize}
  \item \texttt{docs/index.md}
  \item \texttt{docs/kernel/sparse-events.md}
  \item \texttt{docs/kernel/hardware-boundary.md}
  \item \texttt{docs/benchmarks/index.md}
  \item \texttt{docs/reference/status.md}
\end{itemize}

\section*{Sparse-Event Model}

The scheduler model represents system state as a vector
\[
  \mu(t) \in \mathbb{R}^D .
\]

The wake condition is
\[
  \Delta(t) = \|\mu(t)-\mu(t_{last})\|_2,
  \qquad
  \Delta(t) \ge \epsilon(t).
\]

The governor adapts \(\epsilon(t)\) from the observed deviation and elapsed
time. Claims about hardware power, latency, security behavior, or production
boot support require artifacts described in the benchmark and evidence-gate
documentation.

\section*{Claim Boundary}

This note does not claim measured speedup, production binary authentication,
power behavior, or formal verification for the complete repository.
Those claims require current artifacts, baselines, and reproducible commands.

\end{document}
